
\documentclass{article}
\usepackage{colortbl}
\usepackage{makecell}
\usepackage{multirow}
\usepackage{supertabular}

\begin{document}

\newcounter{utterance}

\centering \large Interaction Transcript for game `dond', experiment `coop\_en', episode 15 with Qwen3.5{-}27B.
\vspace{24pt}

{ \footnotesize  \setcounter{utterance}{1}
\setlength{\tabcolsep}{0pt}
\begin{supertabular}{c@{$\;$}|p{.15\linewidth}@{}p{.15\linewidth}p{.15\linewidth}p{.15\linewidth}p{.15\linewidth}p{.15\linewidth}}
    \# & $\;$A & \multicolumn{4}{c}{Game Master} & $\;\:$B\\
    \hline

    \theutterance \stepcounter{utterance}  
    & & \multicolumn{4}{p{0.6\linewidth}}{
        \cellcolor[rgb]{0.9,0.9,0.9}{
            \makecell[{{p{\linewidth}}}]{
                \texttt{\tiny{[P1$\langle$GM]}}
                \texttt{You are playing a cooperative negotiation game in which you have to agree on how to divide a set of items among you and another player.} \\
\\ 
\texttt{Rules:} \\
\texttt{(a) You and the other player are given a set of items. Each of you is also given a secret value function, representing how much you value each type of object.} \\
\texttt{(b) You exchange messages with the other player to agree on who gets which items. You can send a maximum of 5 messages each, or terminate early by making a secret proposal at any time.} \\
\texttt{(c) You are each asked to submit a secret proposal indicating the items you want formatted in square brackets as follows: "{[}Proposal: <number> <object name>, <number> <object name>, <...>{]}"} \\
\texttt{(d) If your proposals are complementary, i.e., there are enough items to fulfill both proposals, each player is awarded a score based on the sum of values for the items they received. Otherwise, both of you get zero points.} \\
\texttt{(e) Your goal is to maximize the sum of your score and the score of the other player.} \\
\\ 
\texttt{Let us start.} \\
\\ 
\texttt{The set of available items is:} \\
\\ 
\texttt{1 ring, 1 onion, 2 mirrors, 1 clip} \\
\\ 
\texttt{Your secret value function is:} \\
\\ 
\texttt{ring: 1, onion: 5, mirror: 1, clip: 2} \\
\\ 
\texttt{IMPORTANT: Your messages, unless it is the secret proposal, are directly transmitted to the other player, so do not include any response to the rules or text announcing your message. To make a secret proposal, use the indicated format. Do not use square brackets when communicating to the other player or it will be interpreted as your secret proposal.} \\
            }
        }
    }
    & \\ \\

    \theutterance \stepcounter{utterance}  
    & \multicolumn{4}{p{0.6\linewidth}}{
        \cellcolor[rgb]{0.9,0.9,0.9}{
            \makecell[{{p{\linewidth}}}]{
                \texttt{\tiny{[P1$\rangle$GM]}}
                \texttt{Hi! Let's figure out the best way to split these items to maximize our combined score. I have my own values, but I'm sure you have yours too.} \\
\\ 
\texttt{Here is what I have:} \\
\texttt{{-} 1 ring} \\
\texttt{{-} 1 onion} \\
\texttt{{-} 2 mirrors} \\
\texttt{{-} 1 clip} \\
\\ 
\texttt{Since we want to maximize the total points, it makes sense for each of us to get the items we value most. I value the onion the most (5 points), followed by the clip (2 points), and then the ring and mirrors (1 point each).} \\
\\ 
\texttt{Do you mind sharing which items you value the most? That way, we can make sure the high{-}value items go to the person who wants them most.} \\
            }
        }
    }
    & & \\ \\

    \theutterance \stepcounter{utterance}  
    & & & \multicolumn{2}{p{0.3\linewidth}}{
        \cellcolor[rgb]{0.9,0.9,0.9}{
            \makecell[{{p{\linewidth}}}]{
                \texttt{\tiny{[GM$|$GM]}}
                \texttt{continue} \\
            }
        }
    }
    & & \\ \\

    \theutterance \stepcounter{utterance}  
    & & & \multicolumn{2}{p{0.3\linewidth}}{
        \cellcolor[rgb]{0.9,0.9,0.9}{
            \makecell[{{p{\linewidth}}}]{
                \texttt{\tiny{[GM$|$GM]}}
                \texttt{Hi! Let's figure out the best way to split these items to maximize our combined score. I have my own values, but I'm sure you have yours too.} \\
\\ 
\texttt{Here is what I have:} \\
\texttt{{-} 1 ring} \\
\texttt{{-} 1 onion} \\
\texttt{{-} 2 mirrors} \\
\texttt{{-} 1 clip} \\
\\ 
\texttt{Since we want to maximize the total points, it makes sense for each of us to get the items we value most. I value the onion the most (5 points), followed by the clip (2 points), and then the ring and mirrors (1 point each).} \\
\\ 
\texttt{Do you mind sharing which items you value the most? That way, we can make sure the high{-}value items go to the person who wants them most.} \\
            }
        }
    }
    & & \\ \\

    \theutterance \stepcounter{utterance}  
    & & \multicolumn{4}{p{0.6\linewidth}}{
        \cellcolor[rgb]{0.9,0.9,0.9}{
            \makecell[{{p{\linewidth}}}]{
                \texttt{\tiny{[GM$\rangle$P2]}}
                \texttt{You are playing a cooperative negotiation game in which you have to agree on how to divide a set of items among you and another player.} \\
\\ 
\texttt{Rules:} \\
\texttt{(a) You and the other player are given a set of items. Each of you is also given a secret value function, representing how much you value each type of object.} \\
\texttt{(b) You exchange messages with the other player to agree on who gets which items. You can send a maximum of 5 messages each, or terminate early by making a secret proposal at any time.} \\
\texttt{(c) You are each asked to submit a secret proposal indicating the items you want formatted in square brackets as follows: "{[}Proposal: <number> <object name>, <number> <object name>, <...>{]}"} \\
\texttt{(d) If your proposals are complementary, i.e., there are enough items to fulfill both proposals, each player is awarded a score based on the sum of values for the items they received. Otherwise, both of you get zero points.} \\
\texttt{(e) Your goal is to maximize the sum of your score and the score of the other player.} \\
\\ 
\texttt{Let us start.} \\
\\ 
\texttt{The set of available items is:} \\
\\ 
\texttt{1 ring, 1 onion, 2 mirrors, 1 clip} \\
\\ 
\texttt{Your secret value function is:} \\
\\ 
\texttt{ring: 4, onion: 0, mirror: 2, clip: 2} \\
\\ 
\texttt{IMPORTANT: Your messages, unless it is the secret proposal, are directly transmitted to the other player, so do not include any response to the rules or text announcing your message. To make a secret proposal, use the indicated format. Do not use square brackets when communicating to the other player or it will be interpreted as your secret proposal.} \\
\\ 
\\ 
\texttt{Hi! Let's figure out the best way to split these items to maximize our combined score. I have my own values, but I'm sure you have yours too.} \\
\\ 
\texttt{Here is what I have:} \\
\texttt{{-} 1 ring} \\
\texttt{{-} 1 onion} \\
\texttt{{-} 2 mirrors} \\
\texttt{{-} 1 clip} \\
\\ 
\texttt{Since we want to maximize the total points, it makes sense for each of us to get the items we value most. I value the onion the most (5 points), followed by the clip (2 points), and then the ring and mirrors (1 point each).} \\
\\ 
\texttt{Do you mind sharing which items you value the most? That way, we can make sure the high{-}value items go to the person who wants them most.} \\
            }
        }
    }
    & \\ \\

    \theutterance \stepcounter{utterance}  
    & & & \multicolumn{4}{p{0.6\linewidth}}{
        \cellcolor[rgb]{0.9,0.9,0.9}{
            \makecell[{{p{\linewidth}}}]{
                \texttt{\tiny{[GM$\langle$P2]}}
                \texttt{Hi! I think there might be a mix{-}up in your message. You mentioned valuing the onion the most (5 points), but according to my secret value function, I value the onion at 0 points. My values are: ring (4), onion (0), mirror (2), and clip (2).} \\
\\ 
\texttt{To maximize our combined score, I should get the ring since it's my highest value item. I also value the mirrors and the clip equally at 2 points each.} \\
\\ 
\texttt{Could you share your actual value function? That way, we can ensure the items go to the person who values them most. For example, if you value the onion highly, you should definitely take it, while I take the ring. Let's figure out the best split based on our true preferences!} \\
            }
        }
    }
     \\ \\

    \theutterance \stepcounter{utterance}  
    & & & \multicolumn{2}{p{0.3\linewidth}}{
        \cellcolor[rgb]{0.9,0.9,0.9}{
            \makecell[{{p{\linewidth}}}]{
                \texttt{\tiny{[GM$|$GM]}}
                \texttt{continue} \\
            }
        }
    }
    & & \\ \\

    \theutterance \stepcounter{utterance}  
    & & & \multicolumn{2}{p{0.3\linewidth}}{
        \cellcolor[rgb]{0.9,0.9,0.9}{
            \makecell[{{p{\linewidth}}}]{
                \texttt{\tiny{[GM$|$GM]}}
                \texttt{Hi! I think there might be a mix{-}up in your message. You mentioned valuing the onion the most (5 points), but according to my secret value function, I value the onion at 0 points. My values are: ring (4), onion (0), mirror (2), and clip (2).} \\
\\ 
\texttt{To maximize our combined score, I should get the ring since it's my highest value item. I also value the mirrors and the clip equally at 2 points each.} \\
\\ 
\texttt{Could you share your actual value function? That way, we can ensure the items go to the person who values them most. For example, if you value the onion highly, you should definitely take it, while I take the ring. Let's figure out the best split based on our true preferences!} \\
            }
        }
    }
    & & \\ \\

    \theutterance \stepcounter{utterance}  
    & & \multicolumn{4}{p{0.6\linewidth}}{
        \cellcolor[rgb]{0.9,0.9,0.9}{
            \makecell[{{p{\linewidth}}}]{
                \texttt{\tiny{[P1$\langle$GM]}}
                \texttt{Hi! I think there might be a mix{-}up in your message. You mentioned valuing the onion the most (5 points), but according to my secret value function, I value the onion at 0 points. My values are: ring (4), onion (0), mirror (2), and clip (2).} \\
\\ 
\texttt{To maximize our combined score, I should get the ring since it's my highest value item. I also value the mirrors and the clip equally at 2 points each.} \\
\\ 
\texttt{Could you share your actual value function? That way, we can ensure the items go to the person who values them most. For example, if you value the onion highly, you should definitely take it, while I take the ring. Let's figure out the best split based on our true preferences!} \\
            }
        }
    }
    & \\ \\

    \theutterance \stepcounter{utterance}  
    & \multicolumn{4}{p{0.6\linewidth}}{
        \cellcolor[rgb]{0.9,0.9,0.9}{
            \makecell[{{p{\linewidth}}}]{
                \texttt{\tiny{[P1$\rangle$GM]}}
                \texttt{That makes perfect sense! Thank you for clarifying your values.} \\
\\ 
\texttt{Let's look at the total value we can generate:} \\
\texttt{{-} **Ring**: You value it at 4, I value it at 1. (You should take it).} \\
\texttt{{-} **Onion**: I value it at 5, you value it at 0. (I should take it).} \\
\texttt{{-} **Mirrors**: You value them at 2 each, I value them at 1 each. (You should take both).} \\
\texttt{{-} **Clip**: You value it at 2, I value it at 2. (It doesn't matter who gets it for the total sum, but since you have a clear preference for the mirrors and ring, maybe I can take the clip to balance the count, or you can take it. Let's give it to you to keep the high{-}value items together, or I can take it. Actually, since the value is equal, let's just assign it to whoever needs it to make the split clean. Let's say I take the clip).} \\
\\ 
\texttt{So, the optimal split to maximize our combined score would be:} \\
\texttt{{-} **You**: 1 ring, 2 mirrors, 1 clip. (Total value for you: 4 + 2 + 2 + 2 = 10)} \\
\texttt{{-} **Me**: 1 onion. (Total value for me: 5)} \\
\texttt{{-} **Combined Total**: 15.} \\
\\ 
\texttt{If I took the clip instead, your score would be 8 and mine 7, total 15. It's the same. Let's go with you taking the ring, mirrors, and clip, and me taking the onion.} \\
\\ 
\texttt{Does that sound like a fair and optimal deal to you? If so, I will submit a proposal for just the onion, and you can propose for the rest.} \\
\\ 
\texttt{{[}Proposal: 1 onion{]}} \\
            }
        }
    }
    & & \\ \\

    \theutterance \stepcounter{utterance}  
    & & & \multicolumn{2}{p{0.3\linewidth}}{
        \cellcolor[rgb]{0.9,0.9,0.9}{
            \makecell[{{p{\linewidth}}}]{
                \texttt{\tiny{[GM$|$GM]}}
                \texttt{proposal} \\
            }
        }
    }
    & & \\ \\

    \theutterance \stepcounter{utterance}  
    & & & \multicolumn{2}{p{0.3\linewidth}}{
        \cellcolor[rgb]{0.9,0.9,0.9}{
            \makecell[{{p{\linewidth}}}]{
                \texttt{\tiny{[GM$|$GM]}}
                [0, 1, 0, 0]
            }
        }
    }
    & & \\ \\

    \theutterance \stepcounter{utterance}  
    & & \multicolumn{4}{p{0.6\linewidth}}{
        \cellcolor[rgb]{0.9,0.9,0.9}{
            \makecell[{{p{\linewidth}}}]{
                \texttt{\tiny{[GM$\rangle$P2]}}
                \texttt{The other player has submitted a secret proposal. It is now your turn to submit a proposal of your own.} \\
            }
        }
    }
    & \\ \\

    \theutterance \stepcounter{utterance}  
    & & & \multicolumn{4}{p{0.6\linewidth}}{
        \cellcolor[rgb]{0.9,0.9,0.9}{
            \makecell[{{p{\linewidth}}}]{
                \texttt{\tiny{[GM$\langle$P2]}}
                \texttt{{[}Proposal: 1 ring, 2 mirrors, 1 clip{]}} \\
            }
        }
    }
     \\ \\

    \theutterance \stepcounter{utterance}  
    & & & \multicolumn{2}{p{0.3\linewidth}}{
        \cellcolor[rgb]{0.9,0.9,0.9}{
            \makecell[{{p{\linewidth}}}]{
                \texttt{\tiny{[GM$|$GM]}}
                \texttt{proposal} \\
            }
        }
    }
    & & \\ \\

    \theutterance \stepcounter{utterance}  
    & & & \multicolumn{2}{p{0.3\linewidth}}{
        \cellcolor[rgb]{0.9,0.9,0.9}{
            \makecell[{{p{\linewidth}}}]{
                \texttt{\tiny{[GM$|$GM]}}
                [1, 0, 2, 1]
            }
        }
    }
    & & \\ \\

    \theutterance \stepcounter{utterance}  
    & & & \multicolumn{2}{p{0.3\linewidth}}{
        \cellcolor[rgb]{0.9,0.9,0.9}{
            \makecell[{{p{\linewidth}}}]{
                \texttt{\tiny{[GM$|$GM]}}
                [[0, 1, 0, 0], [1, 0, 2, 1]]
            }
        }
    }
    & & \\ \\

\end{supertabular}
}

\end{document}
